% ============================================
% 一、选择题 Q11-20
% ============================================
\section{选择题（二）}

% --- Q11: 函数形参 ---
\begin{frame}{第11题}{函数形参}
    \begin{exampleblock}{题目}
        在C语言的函数中，下列正确的说法是 \textbf{【11】} 。
        \begin{itemize}
            \item A. 形参可以是整型，也可以是指针等其它类型
            \item B. 返回类型为void时不能含有形参
            \item C. 全局变量中已声明的标识符名称不能作形参名
            \item D. 必须有形参
        \end{itemize}
    \end{exampleblock}
    \pause
    \begin{alertblock}{解析}
        \footnotesize
        \begin{itemize}
            \item {\color{highlight}A $\checkmark$：形参可以是任何类型}
            \item B $\times$：void函数可以有形参，如 \texttt{void print(int n);}
            \item C $\times$：形参名会屏蔽同名的全局变量（可以同名）
            \item D $\times$：函数可以无形参，如 \texttt{int getchar(void);}
        \end{itemize}
    \end{alertblock}
\end{frame}

% --- Q12: 宏展开 ---
\begin{frame}{第12题}{宏定义展开陷阱}
    \begin{exampleblock}{题目}
        设有宏定义 \texttt{\#define F(r) r+r}，则 \texttt{F(3) * F(3)} 的值为 \textbf{【12】} 。
        \begin{itemize}
            \item A. 36 \quad B. 0 \quad C. 15 \quad D. 表达式有误
        \end{itemize}
    \end{exampleblock}
    \pause
    \begin{alertblock}{解析}
        宏只是文本替换！\\
        \texttt{F(3) * F(3) → 3+3 * 3+3}\\
        乘法优先级高于加法：$3 + 9 + 3 = \textbf{15}$

        \vspace{0.5em}
        {\color{highlight}陷阱：} 很多人误以为是 $(3+3) \times (3+3) = 36$，
        但那需要定义 \texttt{\#define F(r) ((r)+(r))}

        \vspace{0.5em}
        {\color{highlight}答案：C}
    \end{alertblock}
\end{frame}

% --- Q13: 函数嵌套 ---
\begin{frame}{第13题}{函数的定义与调用嵌套}
    \begin{exampleblock}{题目}
        以下正确的描述是 \textbf{【13】} 。
        \begin{itemize}
            \item A. 函数的定义可以嵌套，但函数的调用不可以嵌套
            \item B. 函数的定义不可以嵌套，但函数的调用可以嵌套
            \item C. 函数的定义和函数的调用均可以嵌套
            \item D. 函数的定义和函数的调用均不可以嵌套
        \end{itemize}
    \end{exampleblock}
    \pause
    \begin{alertblock}{解析}
        \begin{itemize}
            \item C语言中\underline{所有函数定义都是平等的}，不能在函数体内定义另一个函数
            \item 但可以在函数中调用其他函数（包括自身→递归）
        \end{itemize}
        {\color{highlight}答案：B（定义不可嵌套，调用可嵌套）}
    \end{alertblock}
\end{frame}

% --- Q14: 指针与数组 ---
\begin{frame}{第14题}{指针下标运算}
    \begin{exampleblock}{题目}
        有下列语句：\texttt{int a[10]=\{1,2,3,4,5,6,7,8\}; int *p = \&a[5];}\\
        则 \texttt{p[-3]} 的值是 \textbf{【14】} 。
        \begin{itemize}
            \item A. 1 \quad B. 2 \quad C. 3 \quad D. 4
        \end{itemize}
    \end{exampleblock}
    \pause
    \begin{alertblock}{解析}
        \footnotesize
        \begin{itemize}
            \item \texttt{p} 指向 \texttt{a[5]}（值为6）
            \item \texttt{p[-3] $\equiv$ *(p-3)}，即向低位偏移3个元素
            \item \texttt{p-3} 指向 \texttt{a[2]}，其值为3
        \end{itemize}
        \begin{center}
            \texttt{a: [1] [2] [3] [4] [5] [6] [7] [8] [?] [?]}\\
            \texttt{\qquad\ \  \^{}p[-3]\qquad\qquad\ \ \^{}p}
        \end{center}
        {\color{highlight}答案：C（3）}
    \end{alertblock}
\end{frame}

% --- Q15: union ---
\begin{frame}{第15题}{联合体（union）大小}
    \begin{exampleblock}{题目}
        \texttt{unionexp\{int i; float j; double k;\} x;}\\
        该联合体在常见计算机上所占字节数为 \textbf{【15】} 。
        \begin{itemize}
            \item A. 4 \quad B. 8 \quad C. 16 \quad D. 232
        \end{itemize}
    \end{exampleblock}
    \pause
    \begin{alertblock}{解析}
        \small
        \begin{itemize}
            \item union 的大小 = \underline{最大成员的大小}
            \item int: 4字节，float: 4字节，double: 8字节
            \item 最大为 double 的 8 字节
        \end{itemize}
        \vspace{0.3em}
        {\color{highlight}对比：} struct 的大小 $\geq$ 各成员大小之和（考虑对齐）

        \vspace{0.3em}
        {\color{highlight}答案：B（8）}
    \end{alertblock}
\end{frame}

% --- Q16: struct 指针 ---
\begin{frame}[fragile]{第16题}{结构体指针成员访问}
    \begin{exampleblock}{题目}
        若有以下程序段，则使用错误的选项是 \textbf{【16】} 。
    \end{exampleblock}
    \begin{lstlisting}[language=C, basicstyle=\footnotesize\ttfamily]
struct student { int num; int age; };
struct student stu[3] = {{1001,20},{1002,19},{1004,20}};
struct student *p;
p = stu;
\end{lstlisting}
    \begin{itemize}
        \item A. \texttt{(*p).num} \quad B. \texttt{(p++)->num} \quad C. \texttt{p = stu.age} \quad D. \texttt{p++}
    \end{itemize}
    \pause
    \begin{alertblock}{解析}
        \small
        \begin{itemize}
            \item A $\checkmark$：解引用后取成员，等价于 \texttt{p->num}
            \item B $\checkmark$：\texttt{p++} 后取成员，合法
            \item {\color{highlight}C $\times$：\texttt{stu}是数组名（数组），不是结构体变量，没有\texttt{.age}成员！}
            \item D $\checkmark$：指针自增，合法
        \end{itemize}
        {\color{highlight}答案：C}
    \end{alertblock}
\end{frame}

% --- Q17: for循环 ---
\begin{frame}{第17题}{for 循环条件分析}
    \begin{exampleblock}{题目}
        \small 设j和k都是int类型，则for循环语句\\
        \texttt{for (j = 0, k = -1; k = 1; j++, k++) printf("*****\textbackslash n");} \textbf{【17】} 。
        \begin{itemize}
            \setlength{\itemsep}{0pt}
            \item A. 循环结束的条件不合法 \quad B. 是无限循环
            \item C. 循环体一次也不执行 \quad D. 循环体只执行一次
        \end{itemize}
    \end{exampleblock}
    \pause
    \begin{alertblock}{解析}
        \footnotesize
        {\color{highlight}关键：} \texttt{k = 1} 是\underline{赋值}，不是判断！

        赋值表达式的值为1（真），所以条件\underline{永远为真}→无限循环。

        \vspace{0.2em}
        {\color{highlight}易错点：} 初学者常把 \texttt{=} 和 \texttt{==} 混淆！

        \vspace{0.2em}
        {\color{highlight}答案：B}
    \end{alertblock}
\end{frame}

% --- Q18: static ---
\begin{frame}{第18题}{存储类别 static}
    \begin{exampleblock}{题目}
        在一个C源文件中，若要定义一个只允许本源文件所有函数使用的全局变量，\\
        则该变量需要使用的存储类别是 \textbf{【18】} 。
        \begin{itemize}
            \item A. auto \quad B. register \quad C. static \quad D. extern
        \end{itemize}
    \end{exampleblock}
    \pause
    \begin{alertblock}{解析}
        \small
        \begin{itemize}
            \item auto：局部变量默认，离开作用域销毁
            \item register：建议编译器放入寄存器（局部变量）
            \item {\color{highlight}C $\checkmark$：static 修饰全局变量→限制作用域为本文件（内部链接）}
            \item extern：声明外部变量，跨文件使用（外部链接）
        \end{itemize}
        \vspace{0.3em}
        {\color{highlight}记忆：} static = "静态+私密"（文件内可见，外部不可见）
    \end{alertblock}
\end{frame}

% --- Q19: 指针类型 ---
\begin{frame}{第19题}{字符串指针数组——类型判断}
    \begin{exampleblock}{题目}
        设有语句：\\
        \texttt{char s[][81] = \{"Student","Teacher","Father","Mother"\}, *ps=s[2];}\\
        题干语句中 \texttt{ps} 的类型是 \textbf{【19】} 。
        \begin{itemize}
            \item A. char \quad B. char * \quad C. char ** \quad D. char[81]
        \end{itemize}
    \end{exampleblock}
    \pause
    \begin{alertblock}{解析}
        \footnotesize
        \begin{itemize}
            \item \texttt{s} 是 \texttt{char[4][81]} 类型
            \item \texttt{s[2]} 是 \texttt{char[81]} 类型（第2行），即 \texttt{char *}
            \item \texttt{*ps = s[2]}：ps 初始化为指向 \texttt{s[2]} 的指针
        \end{itemize}
        {\color{highlight}答案：B（char *）}
    \end{alertblock}
\end{frame}

% --- Q20: 字符串指针取值 ---
\begin{frame}{第20题}{字符串指针输出}
    \begin{exampleblock}{题目}
        在题干基础上，执行语句 \texttt{printf("\%c,\%s,\%c\textbackslash n",*s[1],ps,*ps);} \\
        的输出的结果是 \textbf{【20】} 。
        \begin{itemize}
            \item A. T, Father, F \quad B. Teacher, F, Father
            \item C. Teacher, Father, Father \quad D. 语法错无输出
        \end{itemize}
    \end{exampleblock}
    \pause
    \begin{alertblock}{解析}
        \begin{itemize}
            \item \texttt{*s[1]}：\texttt{s[1]} = \texttt{"Teacher"}，解引用取首字符→\texttt{'T'}
            \item \texttt{ps}：指向 \texttt{s[2]} = \texttt{"Father"}，\texttt{\%s} 输出字符串 →\texttt{Father}
            \item \texttt{*ps}：解引用取 \texttt{'F'}
        \end{itemize}
        输出：\texttt{T,\ Father,\ F}
        \vspace{0.5em}

        {\color{highlight}答案：A}
    \end{alertblock}
\end{frame}
